\documentclass[aspectratio=169]{beamer}

\usepackage[T1]{fontenc}
\usepackage[utf8]{inputenc}
\usepackage[english,russian]{babel}
\usepackage{amsmath}
\usepackage{amssymb}
\usepackage{booktabs}
\usepackage{graphicx}

\usetheme{Madrid}
\usecolortheme{seahorse}
\setbeamertemplate{navigation symbols}{}
\setbeamertemplate{footline}{%
  \leavevmode%
  \hbox{%
    \begin{beamercolorbox}[wd=.5\paperwidth,ht=2.25ex,dp=1ex,left,leftskip=1em]{author in head/foot}%
      \usebeamerfont{author in head/foot}Кузьмин Д.А. --- НИУ ВШЭ, 2026
    \end{beamercolorbox}%
    \begin{beamercolorbox}[wd=.5\paperwidth,ht=2.25ex,dp=1ex,right,rightskip=1em]{title in head/foot}%
      \usebeamerfont{title in head/foot}\insertframenumber{} / \inserttotalframenumber
    \end{beamercolorbox}%
  }%
  \vskip0pt%
}

\title[Детекция дефектов речи]{Детекция дефектов речи с помощью скороговорок}
\author[Кузьмин Д.А.]{Кузьмин Дмитрий Андреевич}
\institute[НИУ ВШЭ]{%
  Факультет компьютерных наук\\
  Программа «Искусственный интеллект»\\[0.4em]
  Научный руководитель: Е.\,О. Кантонистова
}
\date{Москва, 2026}

\begin{document}

% ── Слайд 1. Титульный ──────────────────────────────────────────────────────
\begin{frame}
  \titlepage
\end{frame}

% ── Слайд 2. Мотивация ──────────────────────────────────────────────────────
\begin{frame}{Мотивация}
  \begin{columns}[T]
    \column{0.55\textwidth}
      \textbf{Проблема}
      \begin{itemize}
        \item Нарушения речи у детей влияют на обучение и социальную адаптацию
        \item Диагностика требует экспертной оценки логопеда
        \item Ручная оценка затратна по времени
        \item Регулярный мониторинг динамики труднодостижим
      \end{itemize}
      \vspace{0.8em}
      \textbf{Решение}
      \begin{itemize}
        \item Автоматическая система первичного скрининга
        \item Скороговорки как стандартизированный стимул
        \item ML/DL-подходы для анализа аудиозаписей
      \end{itemize}
    \column{0.42\textwidth}
      \begin{block}{Типы нарушений в работе}
        \small
        Нарушения, связанные с возрастной незрелостью артикуляционной моторики\\[0.5em]
        \textbf{Проверяемые звуки:}\\
        л, р, с, т, ц, ч, ш, щ
      \end{block}
  \end{columns}
\end{frame}

% ── Слайд 3. Постановка задачи ──────────────────────────────────────────────
\begin{frame}{Постановка задачи}
  \begin{block}{Задача бинарной классификации}
    Пусть $X = \{x_i\}_{i=1}^{N}$ — аудиосигналы, $Y = \{y_i\}_{i=1}^{N}$, $y_i \in \{0,1\}$.\\[0.5em]
    Требуется построить отображение
    \[
      f_\theta : \mathcal{X} \rightarrow \{0, 1\}
    \]
    минимизирующее ошибку классификации на отложенной выборке.
  \end{block}

  \vspace{0.8em}
  \begin{columns}[T]
    \column{0.48\textwidth}
      \textbf{Классы:}
      \begin{itemize}
        \item \textbf{good} — нормальное произношение
        \item \textbf{bad} — наличие речевого дефекта
      \end{itemize}
    \column{0.48\textwidth}
      \textbf{Прикладная цель:}
      \begin{itemize}
        \item Telegram-бот для логопедов
        \item Первичный скрининг и мониторинг динамики результатов
      \end{itemize}
  \end{columns}
\end{frame}

% ── Слайд 4. Датасет ────────────────────────────────────────────────────────
\begin{frame}{Описание датасета}
  \begin{columns}[T]
    \column{0.48\textwidth}
      \textbf{Основные параметры:}
      \begin{itemize}
        \item \textbf{Объём:} 2776 аудиозаписей
        \item \textbf{Баланс классов:} ${\sim}2/3$ — \textit{good}, ${\sim}1/3$ — \textit{bad}
        \item \textbf{Контрольные буквы:} л, р, с, т, ц, ч, ш, щ
      \end{itemize}
      \vspace{0.5em}
      \textbf{Длительность аудио (сек.):}
      \begin{itemize}
        \item Мин: 0.99 \quad Макс: 68.66
        \item Средняя: 10.36 \quad Медиана: 8.84
      \end{itemize}
    \column{0.48\textwidth}
      \begin{block}{Ключевые особенности}
        \small
        \begin{itemize}
          \item Каждая запись помечена ${\geq}1$ контрольной буквой
          \item Дисбаланс классов необходимо учитывать при оценке
          \item Ограниченный объём $\Rightarrow$ акцент на предобученные эмбеддинги
        \end{itemize}
      \end{block}
  \end{columns}
\end{frame}

% ── Слайд 5. Метрики ────────────────────────────────────────────────────────
\begin{frame}{Метрики оценки качества}
  \begin{columns}[T]
    \column{0.55\textwidth}
      \[
        \text{Accuracy} = \frac{TP + TN}{TP + TN + FP + FN}
      \]
      \[
        F1_{\text{bad}} = 2 \cdot \frac{\text{Pr}_{\text{bad}} \cdot \text{Re}_{\text{bad}}}{\text{Pr}_{\text{bad}} + \text{Re}_{\text{bad}}}
      \]
      \[
        F1_{\text{macro}} = \frac{F1_{\text{good}} + F1_{\text{bad}}}{2}
      \]
    \column{0.42\textwidth}
      \begin{block}{Иерархия метрик}
        \begin{enumerate}
          \item $\mathbf{F1_{\text{bad}}}$ — ключевая\\
                \small не пропустить дефект
          \item \normalsize $F1_{\text{macro}}$ — баланс по классам
          \item Accuracy — общий индикатор
        \end{enumerate}
      \end{block}
      \vspace{0.5em}
      \small
      При дисбалансе классов \textit{accuracy} одна не достаточна — наивный классификатор даёт ${\sim}67\%$.
  \end{columns}
\end{frame}

% ── Слайд 6. Группа 1: Классические ML ─────────────────────────────────────
\begin{frame}{Группа 1: Классические ML-модели}
  \begin{columns}[T]
    \column{0.48\textwidth}
      \textbf{Признаки:}
      \begin{itemize}
        \item MFCC (мел-частотные кепстральные коэффициенты)
        \item Просодические характеристики (темп, энергия, pitch)
        \item Vocal-tract характеристики
        \item Флаги контрольных букв (вспомогательные)
      \end{itemize}
    \column{0.48\textwidth}
      \textbf{Классификаторы:}
      \begin{itemize}
        \item Logistic Regression
        \item Support Vector Machine (SVM)
        \item Random Forest
      \end{itemize}
      \begin{block}{Роль в работе}
        \small Интерпретируемый baseline и оценка вклада отдельных акустических признаков
      \end{block}
  \end{columns}
\end{frame}

% ── Слайд 7: Группа 2: DL на спектрограммах ────────────────────────────────
\begin{frame}{Группа 2: Нейронные сети на спектрограммах}
  \begin{columns}[T]
    \column{0.48\textwidth}
      \textbf{Входные представления:}
      \begin{itemize}
        \item Лог-мел спектрограммы
        \item Спектрограммы MFCC как 2D-изображения
      \end{itemize}
      \vspace{0.5em}
      \textbf{Архитектуры:}
      \begin{itemize}
        \item 2D CNN (ResNet, EfficientNet, VGG)
        \item CNN + RNN (LSTM, GRU)
        \item CNN + Attention
      \end{itemize}
    \column{0.48\textwidth}
      \begin{block}{Преимущество}
        \small Автоматическое извлечение временно-частотных паттернов без ручного конструирования признаков
      \end{block}
      \vspace{0.5em}
      \begin{block}{Ключевые источники}
        \small CNN на аудио~[Hershey 2017], CNN на детской речи~[Kuo 2022], ResNet+BiLSTM для заикания~[Kourkounakis 2021]
      \end{block}
  \end{columns}
\end{frame}

% ── Слайд 8: Группа 3: Предобученные эмбеддинги ────────────────────────────
\begin{frame}{Группа 3: Предобученные речевые эмбеддинги}
  \begin{columns}[T]
    \column{0.48\textwidth}
      \textbf{Используемые модели:}
      \begin{itemize}
        \item \textbf{HuBERT} — самообучение на неразмеченной речи, фонетические и высокоуровневые паттерны
        \item \textbf{Whisper} — предобучен на ASR, богатые речевые представления
        \item \textbf{VGGish} — предобучен на аудиоклассификации
      \end{itemize}
    \column{0.48\textwidth}
      \textbf{Схемы использования:}
      \begin{itemize}
        \item Заморозка энкодера + лёгкий классификатор
        \item Частичный fine-tuning
        \item Гибридные схемы (HuBERT + LSTM/ResNet)
      \end{itemize}
      \begin{block}{Мотивация}
        \small Ограниченный датасет (2776 записей) — перенос знаний из больших корпусов критически важен
      \end{block}
  \end{columns}
\end{frame}

% ── Слайд 9: Результаты ─────────────────────────────────────────────────────
\begin{frame}{Результаты экспериментов}
  \begin{center}
    \begin{tabular}{lccc}
      \toprule
      \textbf{Группа моделей} & $\mathbf{F1_{\text{bad}}}$ & $\mathbf{F1_{\text{macro}}}$ & \textbf{Accuracy} \\
      \midrule
      Классические ML (baseline) & низкий & низкий & умеренная \\
      DL на спектрограммах       & средний & средний & умеренная \\
      \textbf{Предобученные эмбеддинги} & \textbf{высокий} & \textbf{высокий} & \textbf{высокая} \\
      \bottomrule
    \end{tabular}
  \end{center}
  \vspace{0.8em}
  \begin{columns}[T]
    \column{0.48\textwidth}
      \textbf{Лидеры:}
      \begin{itemize}
        \item HuBERT-эмбеддинги + классификатор
        \item Whisper-эмбеддинги + классификатор
      \end{itemize}
    \column{0.48\textwidth}
      \textbf{Вывод:}\\
      Предобученные модели устойчиво превосходят классические и DL-подходы при ограниченном объёме данных
  \end{columns}
  \vspace{0.5em}
  \small\textit{Примечание: точные значения метрик будут уточнены в финальной версии работы.}
\end{frame}

% ── Слайд 10: Telegram-бот ──────────────────────────────────────────────────
\begin{frame}{Прикладной сервис: Telegram-бот}
  \begin{columns}[T]
    \column{0.55\textwidth}
      \textbf{Целевая аудитория:} логопеды и педагоги\\[0.5em]
      \textbf{Функциональность:}
      \begin{itemize}
        \item Загрузка аудиозаписи чтения скороговорки
        \item Автоматическое предсказание: \textit{good} / \textit{bad}
        \item Мониторинг динамики результатов ребёнка
        \item Быстрая обратная связь без ожидания эксперта
      \end{itemize}
      \vspace{0.5em}
      \textbf{Архитектура:}
      \begin{itemize}
        \item Telegram Bot API
        \item Inference-сервис на лучшей модели (HuBERT/Whisper)
      \end{itemize}
    \column{0.42\textwidth}
      \begin{block}{Ограничения сервиса}
        \small Система не заменяет логопеда — это инструмент первичного скрининга и объективного мониторинга
      \end{block}
  \end{columns}
\end{frame}

% ── Слайд 11: Заключение ────────────────────────────────────────────────────
\begin{frame}{Заключение}
  \begin{columns}[T]
    \column{0.55\textwidth}
      \textbf{Выполнено:}
      \begin{itemize}
        \item Поставлена задача бинарной классификации детской речи
        \item Проведён сравнительный анализ трёх групп подходов
        \item Лучший результат — предобученные эмбеддинги (HuBERT, Whisper)
        \item Реализован Telegram-бот для логопедов
      \end{itemize}

      \vspace{0.5em}
      \textbf{Дальнейшие шаги:}
      \begin{itemize}
        \item Fine-tuning предобученных моделей на расширенном датасете
        \item Мультиклассовая постановка (по типу звука)
        \item Объяснимость: визуализация активаций для логопеда
      \end{itemize}
    \column{0.42\textwidth}
      \begin{block}{Ключевой вывод}
        Предобученные речевые модели (HuBERT, Whisper) обеспечивают наилучшую детекцию дефектов речи при малом объёме размеченных данных
      \end{block}
  \end{columns}
\end{frame}

% ── Слайд 12: Спасибо ───────────────────────────────────────────────────────
\begin{frame}
  \vfill
  \begin{center}
    {\LARGE Спасибо за внимание!}\\[2em]
    {\large Кузьмин Дмитрий Андреевич}\\[0.5em]
    {\normalsize Научный руководитель: Е.\,О. Кантонистова}\\[2em]
    \begin{block}{\centering Тема}
      \centering
      Детекция дефектов речи с помощью скороговорок
    \end{block}
  \end{center}
  \vfill
\end{frame}

\end{document}
